\chapter{Desarrollo asistido por inteligencia artificial}
\label{chap:ia}

\begin{objetivos}
El estudiante podrá usar modelos generativos como colaboradores no confiables, diseñar contexto y criterios de aceptación, operar un modelo local cuando sea viable, evaluar código generado y mantener trazabilidad, privacidad y responsabilidad humana.
\end{objetivos}

\section{Qué cambia y qué permanece}

Un modelo de lenguaje produce continuaciones probables a partir de instrucciones y contexto. Puede resumir, proponer código, generar pruebas, explicar errores y transformar artefactos con gran velocidad. No conoce automáticamente el propósito local, puede inventar APIs, reproducir vulnerabilidades y expresar certeza sin evidencia. Por eso la unidad de trabajo deja de ser ``pedir código'' y pasa a ser ``formular, generar, inspeccionar, ejecutar y decidir''.

Permanece la responsabilidad profesional. Requisitos, arquitectura, seguridad, licencias y aceptación no se delegan al modelo. La fluidez del resultado no es una validación.

\begin{figure}[htbp]
\centering
\begin{tikzpicture}[node distance=1.15cm and 1.0cm, every node/.style={font=\small}]
\node[decision] (obj) {Objetivo y límites};
\node[activity, right=of obj] (ctx) {Contexto mínimo};
\node[activity, right=of ctx] (gen) {Propuesta del modelo};
\node[activity, below=of gen] (rev) {Revisión humana};
\node[evidence, left=of rev] (val) {Pruebas y escáneres};
\node[decision, left=of val] (dec) {Aceptar, corregir o descartar};
\draw[arrow] (obj)--(ctx); \draw[arrow] (ctx)--(gen); \draw[arrow] (gen)--(rev);
\draw[arrow] (rev)--(val); \draw[arrow] (val)--(dec); \draw[arrow] (dec.west) -| (obj.south);
\end{tikzpicture}
\caption{Bucle de asistencia con decisión y evidencia humanas.}
\end{figure}

\section{Anatomía de una solicitud útil}

Una solicitud técnica debe declarar resultado, contexto relevante, restricciones, formato y criterios de aceptación. Incluir todo el repositorio eleva costo y exposición; omitir contratos induce invenciones. Se seleccionan interfaces, errores, convenciones y pruebas cercanas.

\begin{politicaia}
Ejemplo: ``En Python 3.13 y FastAPI, complete únicamente la función \texttt{crear\_reserva}. Respete la interfaz adjunta, no añada dependencias, preserve códigos 201/409 y no registre datos personales. Primero enumere supuestos; luego entregue parche y pruebas pytest. Aceptación: suite existente, casos límite y análisis estático aprobados.''
\end{politicaia}

Pedir que el modelo explique supuestos facilita revisión, pero la explicación también puede ser incorrecta. La verificación decisiva es ejecutar y contrastar con fuentes y contratos. Dividir una tarea limita el radio del error y mejora la posibilidad de atribuir cambios.

\section{Usos y controles}

\begin{table}[htbp]
\centering
\caption{Usos de IA y evidencia mínima.}
\begin{tabularx}{\textwidth}{p{3cm}X X}
\toprule
\textbf{Uso} & \textbf{Riesgo característico} & \textbf{Validación mínima}\\
\midrule
Generar código & API inexistente, caso límite, licencia incierta & Compilar, probar, revisar dependencia y procedencia.\\
Generar pruebas & Repetir el mismo error del código & Derivar casos desde requisitos y mutar una regla.\\
Refactorizar & Cambiar comportamiento silenciosamente & Pruebas caracterizadoras y comparación de contrato.\\
Documentar & Afirmar capacidades no implementadas & Ejecutar cada comando y enlazar evidencia.\\
Revisar seguridad & Falso negativo o consejo obsoleto & Modelo de amenazas, herramientas y revisión especializada.\\
Agente autónomo & Cambios amplios o acciones externas & Entorno aislado, permisos mínimos, límites y aprobación.\\
\bottomrule
\end{tabularx}
\end{table}

\section{Agentes y permisos}

Un agente combina modelo, instrucciones, memoria y herramientas en un ciclo. El riesgo no depende solo del texto que genera, sino de lo que puede leer y modificar. Se aplica mínimo privilegio: repositorio de práctica, secretos ausentes, comandos permitidos, presupuesto de pasos, registro de acciones y aprobación para publicar, desplegar o borrar.

El contenido recuperado desde documentos, incidencias o páginas puede incluir instrucciones maliciosas. Debe tratarse como datos, no como autoridad. Las reglas del curso y del sistema tienen precedencia; una sugerencia encontrada nunca autoriza extraer secretos ni ampliar permisos.

\section{Modelos locales y herramientas libres}

Ollama facilita ejecutar determinados modelos en una estación local; Continue es una interfaz de asistencia de código con código abierto que puede conectarse a modelos locales. ``Local'' mejora control de tránsito de datos, pero no elimina riesgos de licencia, telemetría, vulnerabilidades ni acceso al repositorio. El modelo concreto puede tener una licencia distinta a la herramienta que lo ejecuta: revísela antes de uso institucional o comercial.

\begin{instalacion}
Instale Ollama desde \url{https://ollama.com/download} y compruebe \texttt{ollama --version}. Instale la extensión Continue desde el mercado de su editor o siga \url{https://docs.continue.dev/guides/ollama-guide}. Descargue un modelo de código cuyo tamaño y licencia sean compatibles con el equipo:
\begin{lstlisting}[language=bash]
ollama list
ollama pull qwen2.5-coder:7b
ollama run qwen2.5-coder:7b
\end{lstlisting}
El nombre es un ejemplo reproducible a la fecha de edición, no una recomendación permanente. Verifique catálogo, licencia y hash actuales. Un modelo de 7 mil millones de parámetros cuantizado suele requerir varios GiB de memoria; si el equipo no lo soporta, use el servicio institucional autorizado o realice el ejercicio por parejas sin subir información restringida.
\end{instalacion}

\section{Protocolo de evaluación}

Una comparación válida congela tarea, contexto, versión del modelo, parámetros, entorno y pruebas. Evalúa corrección, seguridad, mantenibilidad, tiempo total y cantidad de intervención humana. Una sola demostración exitosa no permite concluir superioridad.

\begin{validacion}
Prepare cinco tareas pequeñas con pruebas ocultas al modelo. Ejecute una línea base humana y una asistida. Registre versión, solicitud, archivos compartidos, salida, ediciones, pruebas, escaneos, tiempo y resultado. Revise manualmente una muestra de afirmaciones. No use datos personales ni código cuya política prohíba compartir.
\end{validacion}

\section{Riesgos y gobernanza}

El perfil de IA generativa del NIST organiza riesgos y acciones de gestión; la guía de UNESCO enfatiza un enfoque centrado en las personas para educación e investigación. En el aula se aplican mediante transparencia, alfabetización, protección de datos, acceso equitativo y evaluación auténtica.

\begin{politicaia}
Toda entrega declara herramienta y modelo, propósito, fragmentos afectados y validación humana. Se conserva un registro de solicitudes relevantes, no cadenas de pensamiento privadas. Está prohibido ingresar secretos, datos personales o material restringido. El estudiante debe explicar y modificar su solución sin asistencia. La calificación se basa en decisiones y evidencia, no en ocultar o maximizar uso de IA.
\end{politicaia}

Los principales fallos son alucinación, código inseguro, sesgo, exposición de datos, dependencia de proveedor y pérdida de comprensión. Se enfrentan con fuentes, suites independientes, escaneo, minimización de contexto, alternativas interoperables, documentación y defensa oral.

\begin{ejercicio}
\textbf{Laboratorio 11.1 --- experimento controlado.} Con Ollama y Continue, solicite implementar una regla y pruebas de LibreReserva. Compare con una solución sin asistente bajo el protocolo anterior. Entregue registro, diferencias, defectos encontrados y conclusión limitada por la muestra.

\textbf{Laboratorio 11.2 --- agente confinado.} Configure un agente solo sobre una copia del proyecto, sin credenciales ni red externa si la plataforma lo permite. Autorícelo a proponer un refactor, no a confirmar ni desplegar. Revise el parche, ejecute validaciones y catalogue acciones que exigirían aprobación.
\end{ejercicio}

\section{Preguntas de revisión}
\begin{enumerate}
  \item ¿Qué evidencia falsaría la afirmación de que la IA mejoró la productividad?
  \item ¿Qué parte del repositorio no necesita ver el modelo para resolver una tarea concreta?
  \item ¿Cómo se mantiene autonomía intelectual cuando la generación es casi instantánea?
\end{enumerate}
